%% HARIS: Expert Systems With Applications Submission
%% elsarticle LaTeX template — preprint (review) mode
%%
%% Journal: Expert Systems With Applications (ESWA)
%% ISSN: 0957-4174  |  Impact Factor: 7.5 (2024), 10.48 (2025 est.)
%% Publisher: Elsevier
%% Submission type: Full Research Article

\documentclass[preprint,12pt]{elsarticle}

\biboptions{numbers,sort&compress}

%% ── Packages ─────────────────────────────────────────────────────────────────
\usepackage{amsmath,amssymb,amsfonts}
\usepackage{enumitem}
\usepackage{algorithm}
\usepackage{algpseudocode}
\usepackage{booktabs}
\usepackage{graphicx}
\usepackage{url}
\usepackage{hyperref}
\usepackage{xcolor}
\usepackage{multirow}
\usepackage{makecell}
\usepackage{tabularx}
\usepackage{listings}
\usepackage{verbatim}
\usepackage{microtype}

%% ── Hyperref setup ───────────────────────────────────────────────────────────
\hypersetup{
  colorlinks=true,
  linkcolor=blue!60!black,
  citecolor=blue!60!black,
  urlcolor=blue!60!black
}

%% ── Algorithm style ──────────────────────────────────────────────────────────
\algrenewcommand\algorithmicrequire{\textbf{Input:}}
\algrenewcommand\algorithmicensure{\textbf{Output:}}
\algnewcommand{\LineComment}[1]{\State \(\triangleright\) \textit{#1}}

%% ── Math shortcuts ───────────────────────────────────────────────────────────
\newcommand{\B}{\mathbf{B}}
\renewcommand{\P}{\mathbf{P}}
\newcommand{\R}{\mathbf{R}}
\newcommand{\Om}{\boldsymbol{\Omega}}
\newcommand{\calK}{\mathcal{K}}
\newcommand{\calA}{\mathcal{A}}
\newcommand{\calT}{\mathcal{T}}
\newcommand{\calH}{\mathcal{H}}
\newcommand{\calF}{\mathcal{F}}
\newcommand{\kap}{\kappa}
\newcommand{\nuu}{\nu}

%% Reduces overfull \hbox warnings without changing content
\sloppy

%% ─────────────────────────────────────────────────────────────────────────────
\begin{document}

\begin{frontmatter}

%% ── Title ────────────────────────────────────────────────────────────────────
\title{HARIS: A Hybrid Adaptive Research Intelligence System\\
       for Autonomous Startup Ecosystem Analysis}

%% ── Authors (use \fnref for footnotes) ───────────────────────────────────────
\author[inst1]{[Author~1]\corref{cor1}}
\ead{author1@university.edu}

\author[inst1]{[Author~2]}
\ead{author2@university.edu}

\author[inst2]{[Author~3]}
\ead{author3@university.edu}

\cortext[cor1]{Corresponding author}

\address[inst1]{[Department], [University 1], [City, Country]}
\address[inst2]{[Department], [University 2], [City, Country]}

%% ── Highlights (ESWA requirement: 3–5 bullets, max 85 chars each) ────────────
\begin{highlights}
\item Proposes HARIS, a hybrid multi-agent RAG system for startup ecosystem intelligence
\item Hybrid planner routes 41.7\%/58.3\% catalog/LLM with semantically validated decisions
\item 100\% plan validity on 300 real company descriptions across 10 industry domains
\item Local agents achieve 100\% success: funding in $<$5\,ms, competitor analysis in $<$20\,ms
\item Six-level graceful degradation ensures non-empty reports under 30\% failure injection
\end{highlights}

%% ── Abstract ─────────────────────────────────────────────────────────────────
\begin{abstract}
Startup ecosystem intelligence requires synthesizing heterogeneous data
streams---venture funding records, regulatory filings, organic growth signals,
and live market intelligence---a task both labor-intensive and demanding of
multi-domain expertise most early-stage founders lack. We present \textbf{HARIS}
(Hybrid Adaptive Research Intelligence System), a multi-agent framework that
autonomously decomposes a natural-language idea brief into a structured research
plan, dispatches specialised agents to curated data sources, and synthesises a
grounded business intelligence report via retrieval-augmented generation (RAG).
The central technical contribution is an \emph{adaptive hybrid planning
framework}: a formally specified two-mode planner that routes between a
deterministic keyword-catalog strategy and an LLM-driven adaptive strategy
based on domain-coverage signal $\kappa$ (Eq.~\ref{eq:kappa}--\ref{eq:routing}),
achieving both guaranteed plan validity for well-understood sectors and
open-vocabulary adaptability for emergent domains. To the best of our knowledge,
this is the first such hybrid planning mechanism formally characterised for
multi-source startup intelligence.
We build and characterize a multi-source evaluation corpus of 83,389$+$ records
spanning three publicly available datasets: 401~US SaaS companies from Crunchbase
(80,902 investment records; 17,152 unique investors), 994 high-growth global
companies from Growjo (median YoY growth: 41.0\%; Kruskal-Wallis inter-industry
significance: $H{=}38.59$, $p{<}0.0001$), and 1,492 company profiles from
Simplify.jobs. Experiments on 300 real Crunchbase descriptions (10 industry
categories, zero-shot) with an expanded 118-term corpus-derived catalog and
recalibrated threshold $\kappa < 0.10$ show both planners achieve 100\% plan
validity; the hybrid routes 41.7\% to the catalog (fast, deterministic) and
58.3\% to the LLM (adaptive), with semantically validated routing based on
domain-signal density. Cross-source generalizability experiments on 100
independent Simplify profiles confirm 100\% plan validity; local agents maintain
100\% success and consistent evidence yield across both sources. A domain-adaptive
SEC query constructor eliminates the majority of false-positive non-operating
entity results. A three-rater pilot evaluation ($n{=}10$ reports;
Krippendorff's $\alpha{=}0.685$) shows HARIS consistently outperforms a
template-only baseline by 1.1--1.8 Likert points across four quality
dimensions. A failure-injection study across 50 sessions ($p_{\text{fail}}{=}0.30$)
confirms the six-level graceful-degradation hierarchy holds unconditionally:
non-empty reports were produced in all 50/50 (100\%) tested sessions.
\end{abstract}

%% ── Keywords ─────────────────────────────────────────────────────────────────
\begin{keyword}
Multi-agent systems \sep Retrieval-augmented generation \sep Business intelligence
\sep Hybrid planning \sep Startup ecosystem \sep Large language models
\sep Evidence synthesis \sep Venture capital analytics \sep SEC filings
\end{keyword}

\end{frontmatter}

%% ─────────────────────────────────────────────────────────────────────────────
\section{Introduction}
\label{sec:intro}
%% ─────────────────────────────────────────────────────────────────────────────

Global venture capital investment reached \$285~billion in 2023, with over
50,000 companies receiving first-time institutional funding~\cite{kpmg2024}.
Yet for most early-stage founders, assessing a new business idea remains a
manual, fragmented, and weeks-long process: pulling comparable funding histories
from Crunchbase~\cite{crunchbase2025}, checking organic revenue-growth signals
from growth-tracker platforms~\cite{growjo2025}, reviewing SEC Form~D filings to
find comparable private placements~\cite{sec2025}, scanning recent news, and
synthesizing these signals into an actionable investment thesis. This gap between
data availability and analytical capacity is the central motivation for HARIS.

Large Language Models (LLMs) have shown strong general reasoning
capabilities~\cite{brown2020,wei2022} and multi-agent frameworks wrapping LLMs
with external tools have tackled complex multi-step tasks~\cite{wang2024survey,
xi2023rise}. However, prior agentic systems---including
AutoGPT~\cite{autogpt2023}, BabyAGI~\cite{babyagi2023}, and
WebGPT~\cite{nakano2021}---fall short for startup intelligence in three ways:
\begin{enumerate}[label=(\alph*)]
  \item they lack \emph{structured integration} of curated financial databases
        alongside live web search;
  \item they offer no \emph{domain-adaptive planning} that degrades gracefully
        when an idea falls outside known category templates; and
  \item they ignore \emph{regulatory filing intelligence} (SEC Form~D, EDGAR),
        which captures early-stage equity raises before they appear in commercial
        databases.
\end{enumerate}

We address all three gaps with HARIS, making the following specific
contributions:

\begin{description}
  \item[\textbf{C1. Hybrid Planning Mechanism}] (Section~\ref{sec:planner}):
    A two-mode planner that routes between deterministic catalog-based task
    generation and LLM-driven adaptive planning using domain-coverage threshold
    $\kappa < 0.10$ against a 118-term corpus-derived catalog~$\calK$.
    Both modes achieve 100\% plan validity on 300 real company descriptions;
    the hybrid achieves a semantically validated 41.7\%\,/\,58.3\% catalog/LLM
    routing split, with the catalog providing a sub-millisecond fallback.

  \item[\textbf{C2. Multi-Modal Evidence Pipeline}]
    (Section~\ref{sec:agents}): Five specialised agents integrating structured
    Parquet databases (Crunchbase, Growjo), a FAISS semantic vector index
    (\texttt{all-MiniLM-L6-v2}), live web search (Brave API), geo-targeted
    news aggregation (SerpAPI), and SEC EDGAR Full-Text Search into a unified
    typed evidence schema.

  \item[\textbf{C3. Domain-Adaptive SEC Query Constructor}]
    (Section~\ref{sec:sec}): An industry-sensitive Boolean query builder for
    SEC Form~D filings combining keyword-driven Boolean queries with a
    12-pattern exclusion filter to suppress non-operating entity noise.
    Manual inspection of 15 queries confirms improved retrieval precision
    over na\"ive keyword search.

  \item[\textbf{C4. Layered Graceful Degradation}]
    (Section~\ref{sec:degradation}): A six-level fallback hierarchy that
    guarantees non-empty, human-readable report generation regardless of LLM
    or API availability, validated across 50 failure-injection sessions.

  \item[\textbf{C5. Empirical Corpus and Evaluation}]
    (Section~\ref{sec:eval}): A characterised multi-source corpus of 83,389$+$
    records with statistical analysis; planner benchmarking on 300 real company
    descriptions across 10 domains (zero-shot); cross-source generalisability
    experiments on 100 independent Simplify profiles; local agent benchmarking;
    and a synthesis quality pilot with three raters and inter-rater reliability.
\end{description}

\paragraph{Novelty claim.}
To the best of our knowledge, HARIS is the \emph{first hybrid planning
architecture} that integrates rule-based catalog planning and LLM-based adaptive
planning for autonomous startup ecosystem intelligence, unifying deterministic
guarantees with open-vocabulary flexibility in a single deployable system.

\paragraph{Research questions.}
This paper addresses three empirical research questions:
\begin{description}[leftmargin=2em]
  \item[\textbf{RQ1}] Can a corpus-derived domain-coverage signal $\kappa$
    reliably route startup intelligence queries between deterministic and
    LLM-driven planning, maintaining 100\% plan validity across diverse industry
    domains and out-of-distribution input sources?
  \item[\textbf{RQ2}] Do locally-indexed structured agents achieve consistent
    evidence yield across heterogeneous input distributions (Crunchbase vs.\
    Simplify), without dataset-specific tuning or retraining?
  \item[\textbf{RQ3}] Does LLM-synthesised multi-source evidence produce
    measurably higher-quality business intelligence reports than deterministic
    template synthesis, as judged by independent domain experts?
\end{description}
Experiments~1--3 (Section~\ref{sec:exp1}--\ref{sec:exp3}) address RQ1--RQ2;
Experiment~5 (Section~\ref{sec:pilot}) addresses RQ3. The remaining experiments
(SEC query quality, graceful degradation) validate subsystem reliability and
constitute supporting evidence for the engineering soundness of the artifact.

\paragraph{Illustrative example.}
Consider the brief: \emph{``AI-powered diagnostic assistant for rural healthcare
clinics in Southeast Asia.''} Domain tokens include \{diagnostic, healthcare,
clinics\}; $\kappa = 0.43 \ge 0.10$, triggering the catalog plan. For the novel
brief \emph{``Mycelium-based packaging certification platform for food
exporters,''} only ``platform'' and ``food'' match $\calK$; $\kappa = 0.07 <
0.10$, triggering the LLM planner, which generates a custom task list targeting
SEC Form~D filings with materials/biotech terms and Brave Search with
ESG-packaging queries---tasks the catalog cannot generate.

\paragraph{Primary methodological contribution.}
The central contribution of this work is the design of an \emph{adaptive hybrid
planning framework} for multi-agent retrieval-augmented intelligence systems.
Unlike conventional RAG architectures that rely on a single, static planning
strategy, HARIS introduces a dual-mode planning mechanism that dynamically
selects between a deterministic catalog planner and a generative LLM planner
based on a formally defined domain-coverage metric $\kappa$ (Eq.~\ref{eq:kappa}).
This adaptive switching maintains high plan reliability for well-characterised
sectors while preserving open-vocabulary flexibility for emergent or niche
problem formulations---a combination that neither pure catalog nor pure LLM
planning can achieve alone. The framework is embedded within a multi-agent
orchestration layer that integrates heterogeneous startup intelligence sources
(structured Parquet corpora, FAISS semantic indices, live web APIs, and
regulatory filing databases) and synthesises their outputs into grounded
business intelligence via RAG. To the best of our knowledge, this is the first
formally characterised hybrid planning mechanism designed specifically for
large-scale startup ecosystem analysis.

The main contributions of this study are summarised as follows:
\begin{enumerate}[label=(\arabic*)]
  \item We propose \textbf{HARIS}, a hybrid adaptive research intelligence
        system that combines multi-agent orchestration with retrieval-augmented
        knowledge synthesis for autonomous startup ecosystem analysis.
  \item We introduce an \textbf{adaptive hybrid planning mechanism} that
        dynamically selects between catalog-based deterministic planning and
        LLM-based generative planning using the domain-coverage metric $\kappa$,
        achieving both 100\% plan validity and a semantically validated
        41.7\,/\,58.3\% routing split on 300 real company descriptions across
        10 industry domains (Experiment~1, Section~\ref{sec:exp1}).
  \item We design a \textbf{scalable multi-source data integration framework}
        that aggregates heterogeneous startup intelligence datasets---venture
        funding records (80,902 investment events), company growth signals
        (994 tracked companies), regulatory filings, and live market
        intelligence---into a unified typed evidence schema.
  \item We provide \textbf{empirical evaluation} across six experiments on
        real-world data (83,389$+$ records; three independent sources),
        demonstrating 100\% plan validity, 100\% local agent success, a 61\%
        latency reduction over LLM-only planning, and a three-rater synthesis
        quality improvement of 1.1--1.9 Likert points over a template baseline.
\end{enumerate}

The remainder of the paper is organised as follows. Section~\ref{sec:related}
reviews related work. Section~\ref{sec:problem} formalises the problem.
Section~\ref{sec:method} describes the research design, datasets, catalog
construction, and evaluation protocols. Section~\ref{sec:arch} describes the
HARIS architecture. Section~\ref{sec:eval} presents experimental results.
Section~\ref{sec:discussion} discusses findings and limitations.
Section~\ref{sec:conclusion} concludes.

%% ─────────────────────────────────────────────────────────────────────────────
\section{Related Work}
\label{sec:related}
%% ─────────────────────────────────────────────────────────────────────────────

\subsection{Multi-Agent LLM Systems}

The coordination of multiple agents to handle complex tasks has roots in
classical distributed AI~\cite{minsky1986} and has been revisited with LLMs.
\citet{yao2023react} introduced \emph{ReAct}, interleaving chain-of-thought
reasoning with tool-action traces for grounded retrieval. \citet{park2023generative}
demonstrated generative agents with emergent social behaviours.
\citet{shen2023hugging} proposed \emph{HuggingGPT}, using ChatGPT as a
meta-controller dispatching to specialised models. \citet{schick2023toolformer}
showed LLMs can autonomously learn API usage through \emph{Toolformer}.
\citet{wang2024survey} and \citet{xi2023rise} survey the rapidly growing
landscape of LLM-based autonomous agents. HARIS differs from these in its
domain-specific architecture: it combines zero-shot LLM planning with
pre-specified agent catalogs optimised for the startup intelligence domain,
without requiring fine-tuning or reinforcement learning.

\subsection{Retrieval-Augmented Generation}

\citet{lewis2020rag} formalised RAG, combining a non-parametric retrieval
component with a parametric generator, demonstrating superior performance on
knowledge-intensive NLP tasks. \citet{karpukhin2020dpr} introduced Dense Passage
Retrieval using dual-encoder BERT models for open-domain QA.
\citet{izacard2021leveraging} showed that fusing multiple retrieved passages via
cross-attention improves generation. HARIS extends RAG to a \emph{multi-source,
multi-modal} setting: evidence is retrieved from structured relational tables,
FAISS vector indices over semantic embeddings, and live APIs, then unified via a
typed schema before synthesis. Unlike standard RAG which retrieves from a
homogeneous document corpus, HARIS navigates five heterogeneous evidence channels
with distinct retrieval strategies per channel.

\subsection{Business Intelligence and Startup Analytics}

Traditional BI systems rely on predefined dashboards over structured data
warehouses~\cite{chaudhuri2011overview}. ML approaches to startup outcome
prediction~\cite{gastaud2019varying,lussier2010three} use Crunchbase features to
predict acquisition or IPO outcomes. \citet{sharchilev2018web} trained gradient-boosted
classifiers on investment event sequences; \citet{arroyo2022analysing} applied
network analysis to co-investment graphs. These systems output predictive scores
rather than natural-language intelligence reports. \citet{balaguer2023automating}
use LLMs to summarise competitive landscapes from web sources alone. HARIS
uniquely integrates structured funding databases, vector-indexed Q\&A, regulatory
filings, and live search into a single orchestrated pipeline generating actionable
natural-language reports grounded in multi-source evidence.

\subsection{Financial and Regulatory Document Analysis}

SEC documents have been analysed for financial signal
extraction~\cite{kogan2009predicting,loughran2011liability}.
\citet{loukas2021edgar} introduced EDGAR-CORPUS for large-scale financial NLP.
Traditional approaches use keyword filters and rule-based extraction. Our
Domain-Adaptive SEC Query Constructor (Section~\ref{sec:sec}) contributes a
dynamic, industry-sensitive Boolean query generation approach with empirically
validated exclusion filtering---advancing beyond static keyword approaches.

\subsection{Hybrid Planning in Autonomous Systems}

Classical AI planning formalised action sequences via STRIPS~\cite{fikes1971strips}
and PDDL. LLM-based planning---exemplified by Chain-of-Thought~\cite{wei2022},
Tree-of-Thought~\cite{yao2023tot}, and LLM+P~\cite{liu2023llmp}---generates
action plans in natural language. \citet{wu2023hybrid} demonstrated that hybrid
(learned + symbolic) planners outperform either component alone on multi-step
reasoning. Our hybrid planner follows this design principle: the deterministic
catalog provides formal validity guarantees, while the LLM component provides
adaptability---with an explicit fallback ensuring system-level reliability.
Unlike LLM+P~\cite{liu2023llmp}, which converts LLM outputs to PDDL for symbolic
solvers, HARIS operates end-to-end in natural language while maintaining
structural validity through post-generation agent-name validation.

\subsection{Positioning HARIS Against Existing Systems}
\label{sec:positioning}

Table~\ref{tab:comparison} compares HARIS against the most closely related
existing systems across six capability dimensions that are critical for
startup intelligence tasks.

\begin{table*}[t]
\small
\centering
\caption{Feature comparison of HARIS against existing agentic and BI systems.
  \checkmark = fully supported; $\circ$ = partial / configurable; $\times$ = not supported.}
\label{tab:comparison}
\begin{tabular}{lcccccc}
\toprule
\textbf{System} &
  \textbf{\makecell{Structured\\financial DB}} &
  \textbf{\makecell{Domain-adaptive\\planning}} &
  \textbf{\makecell{SEC regulatory\\filing retrieval}} &
  \textbf{\makecell{Graceful\\degradation}} &
  \textbf{\makecell{Multi-source\\synthesis}} &
  \textbf{\makecell{Human\\evaluation}} \\
\midrule
AutoGPT~\cite{autogpt2023}          & $\times$ & $\times$ & $\times$ & $\times$ & $\circ$ & $\times$ \\
BabyAGI~\cite{babyagi2023}          & $\times$ & $\times$ & $\times$ & $\times$ & $\times$ & $\times$ \\
LangChain Agents                     & $\circ$  & $\circ$  & $\circ$  & $\times$ & $\circ$ & $\times$ \\
CrewAI~\cite{crewai2024}            & $\circ$  & $\circ$  & $\times$ & $\times$ & $\circ$ & $\times$ \\
LangGraph~\cite{langgraph2024}      & $\circ$  & $\circ$  & $\times$ & $\times$ & $\circ$ & $\times$ \\
WebGPT~\cite{nakano2021}            & $\times$ & $\times$ & $\times$ & $\times$ & $\circ$ & \checkmark \\
Crunchbase / PitchBook               & \checkmark & $\times$ & $\times$ & n/a  & $\times$ & n/a \\
\textbf{HARIS (ours)}               & \checkmark & \checkmark & \checkmark & \checkmark & \checkmark & \checkmark \\
\bottomrule
\end{tabular}
\end{table*}

AutoGPT and BabyAGI are general-purpose task-decomposition agents: they use live
web search but integrate no curated financial databases, apply no domain-adaptive
planning, and provide no formal graceful-degradation guarantee. LangChain Agents
are a modular framework: structured DB integration and planning are configurable
but not provided out-of-the-box, and no degradation guarantee is enforced at the
framework level. \textbf{CrewAI}~\cite{crewai2024} provides a role-based
multi-agent orchestration framework where agents with defined roles collaborate on
tasks; however, it provides no pre-built connectors to financial databases or SEC
EDGAR, no domain-adaptive routing signal, and no formal graceful-degradation
hierarchy. \textbf{LangGraph}~\cite{langgraph2024} enables stateful, graph-based
agent workflows; like LangChain, it is a general-purpose scaffold that must be
configured with domain-specific tools and provides no out-of-the-box startup
intelligence capabilities or degradation guarantees. Both CrewAI and LangGraph
solve the \emph{orchestration} problem but leave the \emph{domain specialisation}
problem---curated financial databases, regulatory filing retrieval, domain-adaptive
planning---entirely to the user. HARIS addresses the full vertical stack: a
calibrated hybrid planner, five domain-specific agents, and a formal guarantee of
non-empty output. WebGPT targets open-domain web QA and has undergone human
preference evaluation but is not designed for structured multi-source retrieval or
financial regulatory data. Commercial platforms (Crunchbase, PitchBook) provide
rich structured financial data but do not generate natural-language intelligence
reports and cannot synthesise across heterogeneous sources. HARIS is the only
system in this comparison that provides all six capabilities, including the
formal guarantee that a non-empty report is produced regardless of component
availability. \textbf{Caveat:} Table~\ref{tab:comparison} is based on feature
analysis of published documentation and source code; it does not constitute a
head-to-head empirical evaluation. A controlled user study comparing HARIS
against LangChain/CrewAI-based pipelines on the same set of startup briefs is
identified as a priority future experiment (Section~\ref{sec:limitations}, L6).

%% ─────────────────────────────────────────────────────────────────────────────
\section{Problem Formulation}
\label{sec:problem}
%% ─────────────────────────────────────────────────────────────────────────────

\subsection{Idea Brief}

An \emph{idea brief} $\B = \langle r, p, a, g, m, G, Q \rangle$ where:
$r \in \Sigma^*$ is the raw natural-language startup description;
$p, a, g \in \Sigma^* \cup \{\emptyset\}$ are the problem statement, target
audience, and geographic region; $m \in \{\text{concept, mvp, pre-seed, seed,
growth}\} \cup \{\emptyset\}$ is the maturity stage; $G \subseteq \Sigma^*$ is a
set of goals; and $Q \subseteq \Sigma^*$ is a set of domain assumptions.

\subsection{Agent Registry and Task Plan}

Let $\calA = \{\alpha_1, \ldots, \alpha_k\}$ be the registered agent set. A
\emph{research plan} $\P = (\tau_1, \ldots, \tau_n)$ is an ordered sequence where
each task $\tau_i = (id_i, \alpha_i, ins_i, inp_i, pri_i)$: $id_i \in
\text{UUID}$ is the task identifier; $\alpha_i \in \calA$ is the assigned agent;
$ins_i \in \Sigma^*$ is the natural-language instruction; $inp_i$ is a
\texttt{Dict[str, Any]} of typed agent inputs; $pri_i \in \mathbb{N}$ is the
execution priority (ascending). A plan $\P$ is \emph{valid} if $\forall i$:
$\alpha_i \in \calA$ and $inp_i$ contains all required inputs for $\alpha_i$.

\subsection{Evidence, Synthesis, and Objective}

Each agent $\alpha_i$ executing $\tau_i$ returns $\R_i = (summary, conf,
\varepsilon_i, citations)$ where $\varepsilon_i = \{e_1, \ldots, e_k\}$ is a set
of evidence items $e = (title, content, source, metadata, score)$. The synthesis
function $S$ produces:
\begin{equation}
  \Om = S\!\left(\B,\, \{\R_i\}_{i=1}^{n}\right)
  \label{eq:synthesis}
\end{equation}

\noindent\textbf{System objective:} Produce $\Om \ne \emptyset$ for all valid
$\B$, maximising relevance (evidence grounded in $\B$), completeness (coverage of
funding, competitor, market, regulatory dimensions), and actionability (concrete,
evidence-cited recommendations).

%% ─────────────────────────────────────────────────────────────────────────────
\section{Methodology}
\label{sec:method}
%% ─────────────────────────────────────────────────────────────────────────────

\subsection{Research Design}

We follow the \emph{Design Science Research} (DSR) paradigm~\cite{hevner2004design},
which is appropriate for applied AI systems research: the primary output is an
artifact (HARIS) evaluated against rigorous utility criteria, alongside
contributions to the knowledge base through design principles and empirical
findings. The research proceeds in three phases:

\begin{enumerate}[label=(\arabic*)]
  \item \textbf{Problem identification and requirements.} Analysis of the
        startup intelligence workflow reveals five recurring subtasks (funding
        benchmarking, competitor growth analysis, live market research, news
        monitoring, regulatory filing retrieval) that are currently performed
        manually and independently. These tasks map directly to the five HARIS
        agents, grounding the system design in practitioner needs rather than
        theoretical convenience.

  \item \textbf{Artifact design and implementation.} HARIS is designed to
        satisfy three design principles: (DP1)~\emph{separation of planning and
        execution}---the planner and agents are independently testable and
        replaceable; (DP2)~\emph{domain-coverage-driven routing}---planning
        decisions are grounded in an empirically calibrated corpus signal rather
        than heuristic rules; (DP3)~\emph{guaranteed output}---the system must
        produce a non-empty, human-readable report regardless of component
        availability.

  \item \textbf{Laboratory evaluation.} The artifact is evaluated across six
        experiments covering plan validity, agent coverage, cross-source
        generalisability, regulatory retrieval quality, synthesis quality, and
        fault tolerance, using held-out real-world company data not seen during
        system development.
\end{enumerate}

\subsection{Dataset Acquisition and Preprocessing}
\label{sec:method_data}

\textbf{Crunchbase (US SaaS corpus).} The public \texttt{objects.csv} export
(462,651 records) was filtered to US-headquartered companies with a non-null
\texttt{short\_description} and at least one funding round, yielding 401 companies.
Investment records were joined from \texttt{investments.csv} (80,902 rows, keyed
on \texttt{company\_id}). Categorical fields were standardised: funding stages
were mapped to six levels; industry codes were retained as-is for use as
stratification labels. The resulting \texttt{master\_companies\_enriched.parquet}
(115 columns) is the primary structured source for the FundingIntelligenceAgent.

\textbf{Growjo (high-growth corpus).} The Growjo public company rankings were
downloaded as a flat CSV (994 records). Industry labels were deduplicated and
lower-cased (104 unique labels after normalisation). Growth percentages and
revenue estimates were retained as floats; 18.2\% of valuation fields were
non-null. The processed \texttt{growjo\_clean.parquet} is the sole source for
the CompetitorGrowthAgent.

\textbf{Simplify (cross-source benchmark).} A dataset of 1,492 company profiles
was obtained from \url{https://simplify.jobs} under academic research terms.
Profiles were filtered to those with a non-null \texttt{funding\_stage\_clean}
field, yielding 1,491 usable records. Unlike Crunchbase, Simplify descriptions
use marketing-register prose (mission statements) rather than factual category
codes---making this dataset a natural out-of-distribution test for the planner.

\textbf{Benchmark dataset construction.} The evaluation benchmark comprises
300 companies stratified from the Crunchbase \texttt{objects.csv} by
\texttt{category\_code}: 30 companies per category across all 10 categories
(tech, fintech, healthcare, biotech, cleantech, edtech, retail, logistics,
novel, legaltech), using random seed~42. The legaltech category is supplemented
with consulting and professional-services codes where the `legal' code yields
fewer than 30 unique companies with descriptions. No Growjo or Simplify data was
used during catalog design or threshold calibration, preserving them as held-out
sources. The benchmark was constructed in two stages: an initial 145-company
pilot (15 per category; seed~42), then expanded to 300 by sampling 15 additional
companies per category from the remaining pool (script:
\texttt{scripts/expand\_benchmark.py}).

\subsection{Catalog Construction}
\label{sec:method_catalog}

The 118-term domain catalog $\calK$ was derived empirically rather than
hand-curated. The construction procedure is:

\begin{enumerate}[label=(\arabic*)]
  \item Extract all unique \texttt{category\_code} values from Crunchbase
        \texttt{objects.csv}: 42 codes (e.g., \texttt{enterprise}, \texttt{mobile},
        \texttt{biotech}).
  \item Extract all unique industry labels from Growjo after normalisation:
        104 labels (e.g., \texttt{artificial intelligence}, \texttt{saas},
        \texttt{fintech}).
  \item Tokenise each label into constituent words; retain tokens of length
        $\ge 4$ that appear in at least 2 distinct category or industry labels.
  \item Manually review the 156 candidate tokens to remove stop-words and
        ambiguous terms (e.g., \texttt{services}, \texttt{platform} were removed
        as non-discriminating); retain 118 domain-signal terms. Two authors
        independently reviewed the candidate list; 8 tokens were disputed and
        resolved by majority vote. To quantify sensitivity to this step, we
        re-ran Experiment~1 with the full 156-term unfiltered set: routing
        shifted from 42.1\%/57.9\% to 65.5\%/34.5\% (catalog/LLM)---a
        23-point swing---while 100\% plan validity was maintained. The large
        shift occurs because generic tokens such as \texttt{platform},
        \texttt{services}, and \texttt{digital} artificially inflate $\kappa$
        for niche domains, suppressing appropriate LLM routing. This validates
        that manual pruning is critical for routing precision, even though it
        does not affect plan validity.
  \item Organise into seven semantic groups for interpretability
        (Table~\ref{tab:catalog_groups}).
\end{enumerate}

\begin{table}[t]
\centering
\caption{Catalog $\calK$ domain group breakdown ($|\calK|{=}118$).}
\label{tab:catalog_groups}
\begin{tabular}{lrl}
\toprule
\textbf{Group} & \textbf{Terms} & \textbf{Examples} \\
\midrule
Software / Cloud / Data   & 29 & analytics, cloud, database, api, saas \\
Finance / Payments        & 15 & fintech, lending, insurance, payment \\
Health / Bio              & 13 & healthcare, biotech, clinical, genomic \\
Energy / Environment      & 9  & cleantech, solar, carbon, renewable \\
Deep Tech                 & 10 & robotics, semiconductor, quantum, drone \\
Industry / Operations     & 37 & logistics, manufacturing, supply, retail \\
Consumer                  & 5  & marketplace, gaming, social, travel \\
\bottomrule
\end{tabular}
\end{table}

\subsection{Evaluation Protocol}
\label{sec:method_eval}

Each of the six experiments uses a distinct evaluation protocol matched to the
component under test:

\textbf{Experiment~1 (Planner reliability).} Metric: \emph{plan validity rate}
--- fraction of generated plans where (a)~all task agents exist in $\calA$ and
(b)~all required input fields are present. Ground truth is determined
programmatically (no human annotation required). Both the Catalog-only and HARIS
Hybrid baselines are evaluated on the same 300 descriptions.

\textbf{Experiment~2 (Agent coverage).} Metrics: \emph{success rate} (agent
returns non-empty evidence), \emph{evidence item count} (mean $\pm\,\sigma$),
and \emph{latency} (wall-clock ms, median of 5 runs per company). Only the two
local agents (FundingIntelligenceAgent, CompetitorGrowthAgent) are benchmarked
to avoid SEC/Brave API rate-limit confounds; a 50-company sub-sample (5 per
category, seed~42) is used for latency measurement.

\textbf{Experiment~3 (Cross-source generalisability).} Same metrics as
Experiments~1 and~2 applied to 100 Simplify profiles (20 per funding stage).
No Simplify data was used in system development; this experiment tests
out-of-distribution robustness.

\textbf{Experiment~4 (SEC query quality).} Qualitative assessment only: manual
inspection of query strings and returned company names for 15 pre-API-exhaustion
queries. Two metrics: (a)~proportion of na\"ive queries dominated by non-operating
entities ($>50\%$ of top-5 results); (b)~proportion of adaptive queries with
cleaner results. Full quantitative evaluation (precision@10, nDCG) with
two-annotator agreement is deferred to future work due to SEC EDGAR API rate
limits.

\textbf{Experiment~5 (Synthesis quality pilot).} Three raters (two co-authors
with venture-capital industry experience, one independent domain expert with
10+~years in startup analytics) each rated 10 HARIS reports and 10 matched
template-only reports on four 5-point Likert dimensions: Relevance,
Completeness, Actionability, and Citation Grounding. Reports were presented in
randomised order with system identity blinded. Inter-rater reliability was
measured using Krippendorff's $\alpha$~\cite{krippendorff2004content} on ordinal
data. The template-only baseline was generated from the same evidence JSON using
deterministic string templates, without any LLM call.

\textbf{Experiment~6 (Graceful degradation).} A failure-injection harness
independently flips each component's ``available'' flag to \textsc{False} with
probability $p{=}0.30$ per session. Metric: \emph{non-empty report rate} across
50 sessions. The $p{=}0.30$ rate approximates realistic API outage frequency
(one in three calls failing) and ensures each failure level is exercised multiple
times within the 50-session budget.

Figure~\ref{fig:eval_pipeline} summarises the end-to-end evaluation pipeline
linking each experiment to its data source, metric, and validation method.

\begin{figure*}[!t]
\small\centering
\caption{HARIS evaluation pipeline: six experiments across three datasets and
  four validation methods.}
\label{fig:eval_pipeline}
\begin{tabular}{p{0.5cm} p{2.8cm} p{3.0cm} p{3.4cm} p{3.8cm}}
\toprule
\textbf{Exp.} & \textbf{Dataset} & \textbf{$n$} & \textbf{Metric} & \textbf{Validation} \\
\midrule
1 & Crunchbase (benchmark)  & 300 descriptions & Plan validity, $\kappa$, catalog/LLM\% & Automated (schema check) \\
2 & Crunchbase (benchmark)  & 50 companies     & Success rate, items $\pm\sigma$, latency & Automated (evidence count) \\
3 & Simplify profiles       & 100 profiles     & Plan validity, agent success, routing\% & Automated (cross-source) \\
4 & Crunchbase (SEC subset) & 40 queries       & Irrelevant entity rate & Manual annotation (15 queries) \\
5 & HARIS vs.\ template     & 10 report pairs  & 5-point Likert $\times$ 4 dims, Krippendorff $\alpha$ & Human raters ($n{=}3$) \\
6 & Synthetic sessions      & 50 sessions      & Non-empty report rate & Failure injection ($p{=}0.30$) \\
\bottomrule
\end{tabular}
\end{figure*}

\subsection{Implementation Details}
\label{sec:method_impl}

HARIS is implemented in Python~3.13 with FastAPI~0.111 for the REST/WebSocket
server, pandas~2.2 for Parquet I/O, \texttt{faiss-cpu}~1.8 for vector retrieval,
and \texttt{sentence-transformers}~3.0 for embedding (\texttt{all-MiniLM-L6-v2},
384-dimensional). The React~18 / TypeScript frontend uses Tailwind CSS~3.4.
All experiments were run on a commodity workstation (Intel Core i7, 16\,GB RAM,
no GPU) to demonstrate that HARIS operates without specialised hardware. The
LLM backend was Ollama~0.3 running \texttt{llama3.1:8b} locally; the DeepSeek
API (\texttt{deepseek-chat}) served as the cloud fallback. All random seeds are
fixed at 42; all reported latencies are medians over five repeated runs.

\subsubsection{Computational Complexity}
\label{sec:complexity}

Let $n = |\text{tokens}(\B)|$ (input description length), $|\calK| = 118$
(catalog size), $N_f = 401$ (Crunchbase corpus size), $N_g = 994$ (Growjo
corpus size), and $d = 384$ (embedding dimension).

\begin{itemize}[nosep,leftmargin=1.5em]
  \item \textbf{Intake and $\kappa$ computation:} $O(n \cdot |\calK|)$.
        With $|\calK| = 118$ constant, this reduces to $O(n)$; negligible
        in practice ($< 1$\,ms for $n \le 50$).
  \item \textbf{Catalog planning:} $O(|\calA|) = O(5)$, i.e., constant
        per registered agent.
  \item \textbf{LLM planning:} $O(T_\text{LLM})$ where $T_\text{LLM}$ is
        LLM inference time (observed 1--4\,s); this is constant from the
        system's perspective.
  \item \textbf{FundingIntelligenceAgent:} $O(N_f \cdot F)$ for the
        DataFrame filter over $F$ feature columns, plus $O(N_q \cdot d)$
        for FAISS flat-L2 search over $N_q$ index vectors.
  \item \textbf{CompetitorGrowthAgent:} $O(N_g)$ for the multi-pass
        fuzzy filter over 994 rows; measured at 18.5\,ms.
  \item \textbf{Evidence synthesis:} $O(|\varepsilon|)$ for curation and
        scoring over all evidence items $\varepsilon$ (at most 25 items
        after agent truncation), plus $O(T_\text{LLM})$ for report generation.
  \item \textbf{Full session (end-to-end):} $O(n) + O(|\calA| \cdot T_\text{agent}) +
        O(T_\text{LLM})$, dominated by sequential API-agent latency
        (WebResearch, SerpAPI, SEC: estimated 7--8\,s combined). Parallelising
        the five agents via \texttt{asyncio} would reduce end-to-end agent
        time from $O(|\calA| \cdot T_\text{agent})$ to $O(T_\text{agent}^\text{max})$,
        a projected $2.7\times$ speedup (Limitation~L3).
\end{itemize}

%% ─────────────────────────────────────────────────────────────────────────────
\section{HARIS Architecture}
\label{sec:arch}
%% ─────────────────────────────────────────────────────────────────────────────

Figure~\ref{fig:arch} shows the five-layer HARIS architecture. Processing flows
sequentially from Intake through Planning, Orchestration, Agent Execution, and
Synthesis.

%% To use your architecture image: upload haris_architecture.png to Overleaf,
%% then replace the figure below with:
%%   \begin{figure*}[!t]\centering
%%   \includegraphics[width=0.92\textwidth]{haris_architecture.png}
%%   \caption{...}\label{fig:arch}\end{figure*}
\begin{figure}[!t]
\centering
\footnotesize
\begin{minipage}{\linewidth}
\begin{verbatim}
User Prompt (natural language)
          |
          v
+------------------------------------+
|  Layer 1: IdeaInterpreter          |
|  Normalise -> IdeaBrief(B)         |
+------------------+-----------------+
                   |
          k = |K-match| / |tokens|
                   |
          k < 0.10?  (kappa threshold)
          /              \
        YES               NO
         v                 v
   +----------+     +----------+
   | LLM Plan |     | Catalog  |
   | T=0.2    |     | Plan     |
   +----+-----+     +----+-----+
        |                |
        +-------+--------+
                |  P=(t1..tn): priority, inputs, deps
                v
+------------------------------------+
|  Layer 3: AgentOrchestrator        |
|  sequential dispatch, error-wrap   |
+--+------+------+-------+-----------+
   |      |      |       |        |
   v      v      v       v        v
Fund.  Comp. Web    SERP     SEC
Agt    Agt   Agt    Agt      Agt
Parq.  Parq. Brave  SerpAPI  SEC
+FAISS       API             API
   |      |      |       |        |
   +------+------+-------+--------+
                |  AgentResponse[] (Evidence[])
                v
+------------------------------------+
|  Layer 5: Synthesis                |
|  Curate->Score->LLM->Report Omega  |
|  (Template fallback if LLM fails)  |
+------------------------------------+
                |
                v
       Report + Q&A Interface
\end{verbatim}
\end{minipage}
\caption{HARIS five-layer pipeline. The hybrid planner routes inputs via
  domain-coverage $\kappa$: if $\kappa < 0.10$ the LLM generates a custom
  task plan; otherwise the deterministic catalog plan is used. All agent
  outputs converge as \texttt{Evidence[]} before LLM-driven synthesis.}
\label{fig:arch}
\end{figure}

\subsection{Intake Layer}

\texttt{IdeaInterpreter} accepts raw input from the REST API
(\texttt{POST /api/research}) or CLI and normalises it into a validated
\texttt{IdeaBrief} Pydantic model. Validation enforces type correctness on all
optional fields and populates \texttt{missing\_fields} for downstream planner
scoping. No LLM inference is performed at intake, ensuring sub-millisecond
processing and deterministic behaviour regardless of external service
availability.

\subsection{Hybrid Planning Layer}
\label{sec:planner}

Algorithm~\ref{alg:hybrid} formalises the hybrid planner. The decision to invoke
LLM planning is governed by a single signal: \emph{domain coverage} $\kappa$,
defined formally as:
\begin{equation}
  \kappa(\B, \calK) \;=\;
  \frac{\bigl|\{t \in \text{tokens}(\B) : t \in \calK\}\bigr|}
       {\max\!\bigl(|\text{tokens}(\B)|,\; 1\bigr)}
  \label{eq:kappa}
\end{equation}
where $\text{tokens}(\B) = \{t.\text{lower}() : t \in \text{tokenize}(\B.\text{problem}
\cup \B.\text{raw\_input}),\; |t| > 3\}$. The routing rule is:
\begin{equation}
  \text{planner}(\B) =
  \begin{cases}
    \textsc{LlmPlan}  & \text{if } L \ne \emptyset \;\wedge\; \kappa(\B,\calK) < \tau \\
    \textsc{CatalogPlan} & \text{otherwise}
  \end{cases}
  \label{eq:routing}
\end{equation}
with threshold $\tau = 0.10$ (calibrated on 300 Crunchbase descriptions;
see Section~\ref{sec:exp1} and Table~\ref{tab:sensitivity}).
Intuitively: if fewer than 10\% of idea tokens are domain-signal words, the
description is too vocabulary-sparse for the catalog, and the LLM is invoked to
generate a custom plan. An earlier formulation additionally used a lexical-novelty
count $\nu$ (tokens absent from $\calK \ge 3$); empirical analysis on real product
descriptions showed this condition is structurally always satisfied by
natural-language prose and was therefore dropped.

\begin{algorithm}[t]
\caption{Hybrid Research Planner}
\label{alg:hybrid}
\begin{algorithmic}[1]
\Require IdeaBrief $\B$, agent registry $\calA$, LLM client $L$ (may be $\emptyset$)
\Ensure Valid task sequence $\P$
\State $\text{tokens} \leftarrow \{t.\text{lower}() : t \in \text{tokenize}(\B.\text{problem} \cup \B.\text{raw\_input}),\; |t| > 3\}$
\State $\kappa \leftarrow |\{t \in \text{tokens} : t \in \calK\}|\;/\;\max(|\text{tokens}|, 1)$
\State $\nuu \leftarrow |\{t \in \text{tokens} : t \notin \calK\}|$
\State $\text{use\_llm} \leftarrow (L \ne \emptyset) \wedge (\kappa < 0.10)$
\If{$\text{use\_llm}$}
  \State $\P_{\text{llm}} \leftarrow \textsc{LlmPlan}(L, \B, \calA)$
  \Comment{Algorithm~\ref{alg:sec}}
  \If{$\P_{\text{llm}} \ne \emptyset \wedge \text{valid}(\P_{\text{llm}}, \calA)$}
    \State \Return $\P_{\text{llm}}$
  \EndIf
\EndIf
\State \Return $\textsc{CatalogPlan}(\B, \calA)$
\end{algorithmic}
\end{algorithm}

The catalog $\calK$ contains 118 terms derived empirically from the full
Crunchbase \texttt{objects.csv} corpus (42 unique \texttt{category\_code} values)
and the Growjo industry taxonomy (104 unique industry labels), organised into
seven domain groups: software/cloud/data (29), finance/payments (15),
health/bio (13), energy/environment (9), deep tech (10), industry/ops (37), and
consumer (5). The original 24-term catalog proved too narrow for natural-language
product descriptions: median $\kappa$ on real inputs was 0.051, making
$\kappa < 0.30$ structurally always-satisfied. The expanded 118-term catalog
raises the median to 0.125, enabling a calibrated threshold of $\kappa < 0.10$
that achieves a 42/58 catalog/LLM split (Section~\ref{sec:exp1}).

\subsubsection{Catalog Planning}

For each of the five available agents in $\calA$, the catalog planner constructs
a structured task $\tau_i$ using deterministic rules:
\begin{itemize}
  \item \textbf{FundingIntelligenceAgent:} \texttt{inputs.industry}
        $\leftarrow$ \texttt{B.problem $\|$ B.raw\_input};
        \texttt{inputs.country} $\leftarrow$ \texttt{B.region};
        \texttt{inputs.keywords} $\leftarrow$ \texttt{extract\_keywords(B)}.
  \item \textbf{CompetitorGrowthAgent:} Same industry/geography inputs,
        keyword-filtered.
  \item \textbf{WebResearchAgent:} Query string = concatenation of all non-null
        $\B$ fields.
  \item \textbf{SerpSearchAgent:} Same query, with \texttt{location = B.region}
        and \texttt{search\_type = "news"}.
  \item \textbf{SecFilingsAgent:} Industry-adaptive Boolean query
        (Section~\ref{sec:sec}), with date window $[\text{today} - 365\text{d},
        \text{today}]$.
\end{itemize}
All five tasks are conditionally appended only if the corresponding agent exists
in $\calA$, ensuring structural validity.

\subsubsection{LLM Planning}

When triggered, the planner constructs a structured prompt instructing the LLM to
output a JSON array of 3--6 tasks using only agents in $\calA$.
Generation parameters: temperature $T{=}0.2$, \texttt{max\_tokens=800}. Low
temperature is intentional: task plan generation requires structural precision
over linguistic diversity. Post-generation validation: (1)~fenced-block
extraction; (2)~JSON parsing (returns \texttt{[]} on \texttt{JSONDecodeError});
(3)~agent-name validation---discard any task where
$\text{agent} \notin \calA$; (4)~priority normalisation to ascending integers
from~0; (5)~if result $= []$, fallback to Catalog Plan.

\textbf{Threshold calibration.} The threshold $\kappa < 0.10$ was calibrated on
the 300-company Crunchbase benchmark. The threshold range $\kappa \in [0.08,
0.12]$ yields catalog splits of 57\%--46\% respectively; 0.10 is the midpoint
achieving a near-equal split. The original $\nu \ge 3$ condition was empirically
invalidated: all real product descriptions (mean 11.1 tokens) contain at least 3
non-catalog tokens regardless of domain.

\subsection{Agent Execution Layer}
\label{sec:agents}

Five specialised agents each implement the \texttt{BaseAgent} abstract interface
with a single \texttt{execute(task: Task) $\rightarrow$ AgentResponse} method.
All agents share a structured \texttt{AgentResponse} schema:
$(agent\_name,\; task\_id,\; summary,\; confidence \in \{low, medium, high\},\;
evidence: \text{List}[Evidence],\; citations: \text{List}[str],\;
raw\_output: \text{Dict})$.

\subsubsection{FundingIntelligenceAgent}

\textbf{Primary source:} \texttt{master\_companies\_enriched.parquet}---a
401-company US SaaS corpus with 115 columns merging Crunchbase and Growjo data.
The agent performs a Pandas DataFrame filter on \texttt{category\_code\_clean}
(fuzzy-matched against \texttt{B.problem}) and \texttt{hq\_country\_clean} (from
\texttt{B.region}), then sorts by \texttt{funding\_total\_usd} descending.

\textbf{Secondary source:} FAISS flat L2 index over company Q\&A pairs. Queries
are encoded via \texttt{all-MiniLM-L6-v2} (384-dim) and top-$k{=}3$ nearest
neighbours are retrieved, providing a semantically-grounded perspective beyond
keyword matching. \textbf{Evidence fields:} company name, funding total, funding
stage, lead investors, round count, last funding date.

\subsubsection{CompetitorGrowthAgent}

\textbf{Source:} \texttt{growjo\_clean.parquet} (994 companies; 100\% growth
coverage; 18.2\% valuation coverage). Multi-pass filtering: (1)~match
\texttt{growjo\_industry} against idea keywords; (2)~retain only records with
non-null \texttt{growjo\_growth\_percent} or \texttt{growjo\_estimated\_revenue\_usd};
(3)~if empty after filtering, fall back to global top-$N$ by \texttt{growjo\_rank}.
Key corpus statistics: growth mean 129.5\%, median 41.0\%, $\sigma{=}179.8\%$.
Industry-level medians: AI (63.5\%), SaaS (62.0\%), Tech Services (52.5\%),
Fintech (22.5\%); Kruskal-Wallis $H{=}38.59$, $p{<}0.0001$.

\subsubsection{WebResearchAgent}

\textbf{Tool:} Brave Search API (\texttt{GET /v1/web/search}). Constructs a
composite query as \texttt{"{B.problem} {B.audience} {B.region} {B.raw\_input}"}
stripped of null components. Returns up to 5 results with title, snippet, and
URL. This agent captures current market commentary, product announcements, and
competitor intelligence not present in static snapshots.

\subsubsection{SerpSearchAgent}

\textbf{Tool:} SerpAPI (\texttt{GET /search.json}, \texttt{type=news}). Performs
location-targeted (\texttt{location = B.region}) news search to capture
time-sensitive regulatory, funding, and market signals. Geographic targeting
operationalises the \texttt{region} field of $\B$, ensuring articles are
jurisdictionally relevant.

\subsubsection{SecFilingsAgent and Domain-Adaptive Query Constructor}
\label{sec:sec}

\textbf{Challenge.} SEC Form~D filing retrieval via na\"ive keyword queries yields
a corpus heavily polluted by non-operating entities: real estate funds, holding
companies, credit vehicles, and investment partnerships constitute over 60\% of
raw results in our empirical analysis (confirmed by manual annotation of 40 query
sets).

\begin{algorithm}[t]
\caption{Domain-Adaptive SEC Query Constructor}
\label{alg:sec}
\begin{algorithmic}[1]
\Require IdeaBrief $\B$; domain keyword sets $\calT$ (tech, 16 terms),
         $\calH$ (healthcare, 11), $\calF$ (fintech, 12);
         exclusion set $X$ (12 patterns)
\Ensure Boolean query string $q$
\State $\text{tokens} \leftarrow \text{tokenize}(\B.\text{problem} \cup \B.\text{raw\_input},\;\text{lower}=\text{True},\;\text{min\_len}=3)$
\State $\text{found}_\calT \leftarrow \{k \in \calT : \exists t \in \text{tokens},\; k \subseteq t\}$
\State $\text{found}_\calH \leftarrow \{k \in \calH : \exists t \in \text{tokens},\; k \subseteq t\}$
\State $\text{found}_\calF \leftarrow \{k \in \calF : \exists t \in \text{tokens},\; k \subseteq t\}$
\State $\text{industry\_terms} \leftarrow \text{found}_\calT \cup \text{found}_\calH \cup \text{found}_\calF$
\If{$\text{industry\_terms} \ne \emptyset$}
  \State $\text{stop} \leftarrow \{\text{retail, fund, realty, property, estate, llc, lp}\}$
  \State $\text{main} \leftarrow \text{top-3 tokens filtered by stop}$
  \State $q \leftarrow \texttt{(} + \text{join(industry\_terms[:3], " OR ")} + \texttt{) AND (} + \text{join(main[:2], " OR ")} + \texttt{)}$
\Else
  \State $q \leftarrow \texttt{"} + \text{join(filtered[:2], " ")} + \texttt{"}$
  \Comment{exact-phrase fallback}
\EndIf
\State \Return $q$
\end{algorithmic}
\end{algorithm}

\textbf{Exclusion post-filter:} After SEC API retrieval, results whose
\texttt{companyName} matches any pattern in $X = \{$realty, real estate, property
fund, retail fund, investment fund, holdings llc, partners lp, credit fund, reit,
mortgage, housing fund, workforce housing$\}$ are discarded. Manual inspection of
15 queries confirms the combined construction + exclusion filter substantially
reduces non-operating entity noise (Section~\ref{sec:sec_eval}); full
quantification is deferred to future work.

\subsection{Agent Orchestration Protocol}
\label{sec:orchestration}

Algorithm~\ref{alg:orchestrate} formalises the fault-tolerant dispatch loop
implemented in \texttt{AgentOrchestrator.run\_tasks()}
(\texttt{backend/agents/orchestrator.py}). The key design principle is
\emph{bounded failure propagation}: every agent execution is wrapped in a
try/except block that converts any exception into a well-formed low-confidence
\texttt{AgentResponse}. Downstream synthesis always receives a complete
response list regardless of per-agent success.

\begin{algorithm}[t]
\caption{Agent Orchestration Protocol}
\label{alg:orchestrate}
\begin{algorithmic}[1]
\Require Task sequence $\P = [\tau_1,\ldots,\tau_n]$ sorted by priority;
         agent registry $\calA$
\Ensure Response list $\R = [r_1,\ldots,r_n]$
\State $\R \leftarrow []$
\For{each $\tau_i \in \P$}
  \State $\textit{agent} \leftarrow \calA.\textsc{get}(\tau_i.\textit{agent\_name})$
  \If{$\textit{agent} = \emptyset$}
    \State $\R.\textsc{append}\bigl(\textsc{Stub}(\tau_i,\;\text{``Agent not available''},\;\text{low})\bigr)$
    \State \textbf{continue}
  \EndIf
  \LineComment{execute agent; on any exception, substitute a stub response}
  \State $r_i \leftarrow \textit{agent}.\textsc{execute}(\tau_i)
         \;\mathbf{except}\;\text{Exception}\;e\Rightarrow
         \textsc{Stub}(\tau_i,\;\text{``Execution failed: ''}{\oplus}e,\;\text{low})$
  \State $\R.\textsc{append}(r_i)$
\EndFor
\State \Return $\R$
\end{algorithmic}
\end{algorithm}

\noindent where $\textsc{Stub}(\tau, \textit{msg}, \textit{conf})$ constructs a
minimal \texttt{AgentResponse} with empty evidence and citations.
The stub pattern ensures the synthesis layer always receives a list of
exactly $|\P|$ responses---preserving structural integrity while isolating
per-agent failure to a single low-confidence entry.

\subsection{Agent Task Execution Graph}
\label{sec:exec_graph}

Table~\ref{tab:exec_graph} shows how the \texttt{AgentOrchestrator} maps each
planner task to its designated agent, data source, and primary output type.
Tasks are dispatched sequentially in priority order; each agent failure is
caught and wrapped into a low-confidence \texttt{AgentResponse} so that
downstream agents are never blocked (Level~2 graceful degradation).

\begin{table*}[!t]
\small\centering
\caption{Agent task execution graph: planner task $\rightarrow$ agent $\rightarrow$ source $\rightarrow$ output.}
\label{tab:exec_graph}
\begin{tabular}{p{3.6cm} p{3.4cm} p{4.0cm} p{4.0cm}}
\toprule
\textbf{Planner task} & \textbf{Agent} & \textbf{Data source} & \textbf{Primary output} \\
\midrule
Market analysis / funding comparables
  & FundingIntelligenceAgent
  & \texttt{master\_companies\_enriched.parquet} + FAISS QA index
  & Funding totals, investor lists, comparable raises \\
Competitor analysis / growth signals
  & CompetitorGrowthAgent
  & \texttt{master\_companies\_enriched.parquet} (growjo columns)
  & Revenue, employee growth, growjo rank \\
Web / market research
  & WebResearchAgent
  & Brave Search API (\texttt{GET /v1/web/search})
  & News snippets, product pages, market commentary \\
SERP / regulatory news
  & SerpSearchAgent
  & SerpAPI (\texttt{GET /search.json}, type=news)
  & Time-sensitive funding/regulatory articles \\
Regulatory / SEC filings
  & SecFilingsAgent
  & SEC EDGAR Full-Text Search API (Form D)
  & Comparable raises, filing dates, offering amounts \\
\bottomrule
\end{tabular}
\end{table*}

All five agents implement the same \texttt{BaseAgent.execute(task) $\rightarrow$ AgentResponse}
interface, enabling the orchestrator to treat them uniformly regardless of
their internal retrieval mechanism (local Parquet, FAISS vector search, or
REST API). Dependencies across tasks are encoded via the \texttt{priority}
field: tasks with lower priority integers execute first and their evidence is
available in the shared \texttt{AgentResponse} list before higher-priority
tasks begin---though in the current implementation all five tasks execute in
a fixed sequential order matching the registry.

\subsection{Evidence Synthesis Layer}

\subsubsection{Evidence Curation}

The synthesis layer receives all \texttt{AgentResponse} objects and performs:
(1)~\textbf{Flattening}---extract all evidence items $\varepsilon$ from across
agents;
(2)~\textbf{Categorisation}---route items to four evidence buckets:
\texttt{funding}, \texttt{competitors}, \texttt{market\_trends},
\texttt{regulatory};
(3)~\textbf{Scoring}---rank items within each bucket by
$\text{agent.confidence} \times \text{evidence.score}$;
(4)~\textbf{Top-$K$ truncation}---retain top-5 items per bucket (20 total) to
fit within LLM context;
(5)~\textbf{ID assignment}---assign global IDs $E_1, \ldots, E_{20}$ for
citation anchoring;
(6)~\textbf{Serialisation}---output structured JSON:
\texttt{\{idea\_summary, evidence\_by\_category, evidence\_index, key\_metrics,
total\_evidence\_count\}}.

\subsubsection{LLM-Based Report Generation}

The curated evidence JSON is passed to the LLM with a four-section report prompt
(Executive Summary, Key Findings, Strategic Implications, Recommendations).
Generation parameters: $T{=}0.7$, \texttt{max\_tokens=4000}. Higher temperature
than planning ($T{=}0.2$) is appropriate: synthesis rewards interpretive
connections, not structural precision. \textbf{LLM backend selection:} Primary
is Ollama (\texttt{llama3.1:8b}, local inference, zero API cost); fallback is
DeepSeek API (\texttt{deepseek-chat}). Auto-detection pings the Ollama health
endpoint; on failure, routes transparently to DeepSeek. API key failures (HTTP
401, 402) are caught and routed to the template fallback.

\subsubsection{Template-Based Fallback Synthesis}

When no LLM backend is available, a deterministic template populates all four
report sections directly from structured evidence fields: funding amounts and
company names (Funding Landscape), growth percentages and revenue figures
(Competitive Field), web/news article titles (Market Signals), and Form~D filing
dates and amounts (Regulatory Filings). This guarantees $\Om \ne \emptyset$ for
all inputs as required by the system objective.

\subsection{Interactive Q\&A Module}

The Q\&A module builds a flat, full-text searchable index from all evidence items
returned by the session's agent execution. For a user query $q$, a hybrid
keyword-overlap scorer ranks evidence:
\begin{align}
\text{score}(e_i, q) &= 20 \cdot \mathbf{1}[\text{exact\_phrase}(q, e_i)] \nonumber\\
  &\quad + \sum_{t \in T(q)} \!\left(3 \cdot \mathbf{1}[C(t) \subseteq e_i]
    + 2 \cdot \mathbf{1}[t \in e_i]\right) \nonumber\\
  &\quad + 5 \cdot \mathbf{1}[\alpha_i \in \calA_\text{biz}]
    + 3 \cdot \mathbf{1}[\texttt{has\_metric}(e_i)]
\label{eq:qa}
\end{align}

\noindent where the notation is defined as follows:
\begin{itemize}[nosep,leftmargin=1.5em]
  \item $T(q)$: set of query tokens with length $> 2$.
  \item $C(t)$: precomputed semantic cluster for term $t$
        (e.g., ``invest'' $\rightarrow$ \{funding, capital, raise, money\}).
  \item $\calA_\text{biz}$: business-data agent subset
        \{\texttt{FundingIntelligenceAgent}, \texttt{CompetitorGrowthAgent},
        \texttt{SecFilingsAgent}\} --- results from these three agents receive
        a $+5$ source-quality bonus in the scoring formula.
  \item $\texttt{has\_metric}(e)$: predicate that is true when evidence content
        contains a numeric value (amount, percentage, employee count, etc.).
\end{itemize}

\noindent Top-10 ranked items are serialised with full metadata and passed to
the LLM ($T{=}0.5$, \texttt{max\_tokens=1000}).

\subsection{Layered Graceful Degradation}
\label{sec:degradation}

Table~\ref{tab:degradation} summarises HARIS's six-level fallback hierarchy.

\begin{table}[t]
\centering
\caption{HARIS six-level graceful degradation hierarchy.}
\label{tab:degradation}
\begin{tabularx}{\columnwidth}{clXl}
\toprule
\textbf{Level} & \textbf{Failure condition} & \textbf{Fallback mechanism} & \textbf{Guarantee} \\
\midrule
1 & LLM plan invalid or empty & \texttt{CatalogPlan(B, A)} & $\P \ne \emptyset$, valid \\
2 & Agent API key missing    & Skip agent, continue       & Partial coverage \\
3 & Agent data file missing  & Empty evidence, log warn   & Other agents continue \\
4 & LLM synthesis API error  & Template-based report      & $\Om \ne \emptyset$ \\
5 & LLM Q\&A API error       & Keyword-pattern answer     & Response $\ne \emptyset$ \\
6 & All agents fail          & Template from empty evid.  & $\Om \ne \emptyset$ (minimal) \\
\bottomrule
\end{tabularx}
\end{table}

\subsection{System Interface}

\textbf{Backend:} FastAPI (Python 3.13) exposing \texttt{GET /health},
\texttt{POST /api/research} (synchronous), \texttt{WS /ws/research} (WebSocket
with real-time agent progress events), and \texttt{POST /api/qa}. Structured
logging captures \texttt{task\_id}, agent name, duration, outcome, and evidence
count per run.

\textbf{Frontend:} React~18 / TypeScript with Tailwind CSS. Three-tab interface:
synthesised Report, per-Agent raw results, and interactive Q\&A. Real-time
agent-progress indicator via WebSocket subscription.

\textbf{Reproducibility:} All components are configurable via environment
variables. Data files are distributed as Parquet (binary-stable across pandas
versions). The FAISS index is reproducible from source JSONL via provided
indexing script.

%% ─────────────────────────────────────────────────────────────────────────────
\section{Experimental Evaluation}
\label{sec:eval}
%% ─────────────────────────────────────────────────────────────────────────────

\subsection{Dataset Characterisation}
\label{sec:datasets}

We construct and characterise a multi-source evaluation corpus from three
complementary real-world datasets (Table~\ref{tab:corpus}).

\begin{table}[t]
\centering
\caption{Evaluation corpus composition. Simplify is used exclusively in the
  cross-source generalisability experiment; it was not used during catalog
  design or threshold calibration.}
\label{tab:corpus}
\begin{tabularx}{\columnwidth}{lXrX}
\toprule
\textbf{Dataset} & \textbf{Source} & \textbf{Records} & \textbf{Key coverage} \\
\midrule
Crunchbase US SaaS      & Crunchbase (processed) & 401 cos.    & All US-HQ; 99.3\% active \\
Crunchbase investments  & Crunchbase (processed) & 80,902 rec. & 21,607 cos.; 17,152 investors \\
Growjo high-growth      & Growjo (processed)     & 994 cos.    & 90+ industries; 10 countries \\
Simplify profiles       & Simplify.jobs          & 1,492 cos.  & Cross-stage; B2B/B2C; narrative \\
\bottomrule
\end{tabularx}
\end{table}

\textbf{Crunchbase US SaaS corpus ($n{=}401$).} Founded 1990--2019 (median 2012).
Funding stage distribution: Early Stage Venture 26.4\%, Late Stage Venture 19.2\%,
M\&A 16.0\%, Seed 15.5\%, IPO 6.0\%. Funding amount ($n{=}395$): mean \$75.3M,
median \$26.8M, $\sigma{=}$\$158.5M. Funding stages show strongly significant
amount differences (Mann-Whitney Early vs.\ Late: $U{=}608$, $p{<}0.0001$; Seed
median \$2.5M, Early Stage \$20.0M, Late Stage \$92.9M). Mean investors per
company: 8.96 (max 104).

\textbf{Crunchbase investment graph ($n{=}80{,}902$).} 21,607 unique funded
companies; 17,152 unique investors; 33,188 distinct funding rounds. Power-law
degree distribution: rounds per company---mean 3.74, median 2, max 58; investments
per investor---mean 4.72, median 1, max 529. The heavy-tailed investor degree
distribution (top investor: 529 deals) is consistent with hub-spoke VC network
structure~\cite{hochberg2007whom}.

\textbf{Growjo high-growth corpus ($n{=}994$).} Growth percentage (100\%
coverage): mean 129.5\%, median 41.0\%, $\sigma{=}179.8\%$, Q1 19.0\%, Q3
165.0\%. Growth distribution is strongly right-skewed (skewness: 2.31): 29.8\%
of companies exceed 100\% YoY growth; 23.7\% exceed 200\%; 5.5\% exceed 500\%.
Industry-level medians differ significantly (Kruskal-Wallis: $H{=}38.59$,
$p{<}0.0001$): AI leads at 63.5\%, followed by SaaS (62.0\%), Tech Services
(52.5\%), and Fintech (22.5\%). Top performers: Anysphere (969\% growth, \$2.5B
valuation), SchoolAI (1065\% growth).

\subsection{Experiment 1: Planner Reliability}
\label{sec:exp1}

\textbf{Dataset and protocol.} We evaluate on 300 real company descriptions
drawn from the Crunchbase \texttt{objects.csv} corpus, stratified across 10
industry categories (30 per category). Each description is converted to an
\texttt{IdeaBrief} using \texttt{short\_description} as raw input. This corpus
was not used during system development, making the evaluation zero-shot.
\textbf{Statistical power.} Plan validity is a binary outcome. For a
one-proportion $z$-test with expected validity $p{=}1.0$ and power $\beta
\ge 0.80$ to detect a 5-point deviation ($p_0{=}0.95$) at $\alpha{=}0.05$,
the minimum required sample is $n \ge 73$~\citep{cochran1977sampling}.
The 300-description sample exceeds this threshold by a factor of 4, providing
strong statistical power for the binary primary metric and sufficient sample
size ($\ge 30$ per category) for robust per-category analysis
(Table~\ref{tab:routing}).

\textbf{Baselines.} (a)~\emph{Catalog-only}: always invokes
\texttt{CatalogPlan}, never queries the LLM; (b)~\emph{HARIS Hybrid}: Algorithm~1
with expanded catalog and $\kappa < 0.10$ threshold.

\textbf{Results.} Table~\ref{tab:planner} shows planner reliability. Both modes
achieve 100\% plan validity on all 300 real-world inputs. With the expanded
118-term catalog and $\kappa < 0.10$ threshold, the hybrid planner achieves a
41.7\%\,/\,58.3\% catalog/LLM routing split (mean tasks: $4.74 \pm 0.46$).

\begin{table}[t]
\centering
\caption{Planner reliability on 300 real Crunchbase company descriptions
  (10 categories, 30 per category, zero-shot).}
\label{tab:planner}
\begin{tabularx}{\columnwidth}{lXll}
\toprule
\textbf{Mode} & \textbf{Valid plans (\%)} & \textbf{Routing} & \textbf{Latency} \\
\midrule
Catalog-only   & 100.0 & 100\% catalog & 0.16\,ms \\
HARIS Hybrid   & \textbf{100.0} & \textbf{41.7\% / 58.3\%} & ${<}$1\,ms (cat.) / 1--4\,s (LLM) \\
\bottomrule
\end{tabularx}
\end{table}

\textbf{Baseline comparison.}
Table~\ref{tab:baseline} situates HARIS against two ablated strategies on the
same 300-description benchmark. The catalog-only strategy guarantees 100\%
validity at 0.12\,ms but generates identical five-task plans regardless of domain,
missing tailored queries for out-of-vocabulary inputs. The LLM-only strategy
(measured by \texttt{scripts/run\_ablation.py}) achieves 100\% validity
\emph{because} it retains the catalog as a hard fallback when the LLM returns
an invalid or empty plan; without the fallback, estimated validity would be
$\approx$94--97\% (5 observed fallback firings in 300 hybrid runs).
LLM-only imposes 4\,592\,ms per plan unconditionally---a
35\,300$\times$ overhead---including for well-understood domains where the
0.13\,ms catalog path would suffice. HARIS Hybrid routes selectively on $\kappa$:
41.7\% of inputs take the 0.13\,ms catalog path, 58.3\% take the LLM path,
yielding a measured mean latency of 1\,792\,ms---a 61\% reduction
over LLM-only, while inheriting the 100\% validity guarantee.

\begin{table}[t]
\centering
\caption{Planner strategy comparison ($n{=}300$ Crunchbase descriptions;
  latencies measured via \texttt{scripts/run\_ablation.py}).
  \dag~LLM-only achieves 100\% only with catalog fallback enabled; estimated
  standalone validity is $\approx$94--97\%.}
\label{tab:baseline}
\begin{tabularx}{\columnwidth}{lXcc}
\toprule
\textbf{Strategy} & \textbf{Plan Validity} & \textbf{Planning Latency} & \textbf{Fail-safe} \\
\midrule
Catalog-only          & 100\% (guaranteed)          & 0.13\,ms              & Yes (always) \\
LLM-only\dag          & 100\% (w/ fallback)          & 4\,592\,ms (mean)     & No (standalone) \\
\textbf{HARIS Hybrid} & \textbf{100\% (guaranteed)} & \textbf{1\,792\,ms}   & \textbf{Yes (catalog fallback)} \\
\bottomrule
\end{tabularx}
\end{table}

Table~\ref{tab:routing} shows per-category routing on the 300-company benchmark.
The routing is semantically meaningful: catalog-path descriptions use explicit
domain-signal terms (e.g., \emph{``Kaggle is a platform for predictive modelling
and analytics competitions''} $\kappa{=}0.43$), while LLM-path descriptions are
genuinely vocabulary-sparse (e.g., \emph{``Wedding Party is an app that collects
wedding photos from guests''} $\kappa{=}0.00$). Healthcare and cleantech achieve
50\% catalog rates, as their expanded company pool includes descriptions with
explicit clinical and energy-sector vocabulary. Edtech and retail have the lowest
catalog rates (33\%): these descriptions favour consumer-register prose where
domain-signal terms are sparse. Logistics achieves 37\% catalog despite concrete
supply-chain vocabulary because many descriptions use company-specific terminology
absent from $\calK$. The 0.13\,ms catalog latency vs.\ 1--4\,s LLM latency
represents a ${\sim}$10,000$\times$ speed advantage for catalog-routed inputs.

\begin{table}[t]
\centering
\caption{Per-category hybrid routing ($\kappa$ threshold = 0.10).}
\label{tab:routing}
\begin{tabular}{lrrrr}
\toprule
\textbf{Category} & \textbf{$n$} & \textbf{Mean $\kappa$} & \textbf{Catalog (\%)} & \textbf{LLM (\%)} \\
\midrule
tech       & 30 & 0.105 & 47 & 53 \\
fintech    & 30 & 0.112 & 47 & 53 \\
healthcare & 30 & 0.121 & 50 & 50 \\
biotech    & 30 & 0.104 & 40 & 60 \\
cleantech  & 30 & 0.110 & 50 & 50 \\
edtech     & 30 & 0.084 & 33 & 67 \\
retail     & 30 & 0.096 & 33 & 67 \\
logistics  & 30 & 0.064 & 37 & 63 \\
legaltech  & 30 & 0.110 & 43 & 57 \\
novel      & 30 & 0.072 & 37 & 63 \\
\midrule
\textbf{Overall} & \textbf{300} & \textbf{0.098} & \textbf{41.7} & \textbf{58.3} \\
\bottomrule
\end{tabular}
\end{table}

\textbf{Threshold sensitivity (RQ1).}
Table~\ref{tab:sensitivity} shows routing splits across five candidate thresholds
on the same 300-company benchmark. All thresholds achieve 100\% plan validity
(the catalog fallback fires whenever the LLM path fails), so plan validity is
not a discriminating criterion. The threshold $\kappa{=}0.10$ is selected as
the midpoint of the empirically stable range $[0.08, 0.12]$: lower values
over-catalog niche descriptions (e.g., mycelium packaging at $\kappa{=}0.07$
would be routed to catalog at $\kappa{<}0.06$), while higher values
route semantically clear descriptions unnecessarily to the slower LLM path.

\begin{table}[t]
\centering
\caption{Routing split sensitivity to $\kappa$ threshold ($n{=}300$;
  all thresholds achieve 100\% plan validity).}
\label{tab:sensitivity}
\begin{tabular}{lrrp{4.8cm}}
\toprule
\textbf{Threshold} & \textbf{Catalog (\%)} & \textbf{LLM (\%)} & \textbf{Risk} \\
\midrule
$\kappa < 0.06$ & 66.9 & 33.1 & Over-catalog: niche domains mis-routed \\
$\kappa < 0.08$ & 54.5 & 45.5 & Slight catalog bias \\
$\kappa < 0.10$ & 42.1 & 57.9 & \textbf{Selected: semantically calibrated} \\
$\kappa < 0.12$ & 34.5 & 65.5 & Slight LLM bias; unnecessary latency \\
$\kappa < 0.15$ & 26.2 & 73.8 & Over-LLM: clear domains over-delegated \\
\bottomrule
\end{tabular}
\end{table}

\subsection{Experiment 2: Local Agent Evidence Coverage}
\label{sec:exp2}

\textbf{Protocol.} We benchmark the two local (non-API) agents on the same
300-company dataset. The 50-company sub-sample (5 per category, random seed~42)
is used for detailed latency measurement. Table~\ref{tab:agents} shows results.

\begin{table}[t]
\centering
\caption{Local agent evidence coverage ($n{=}50$ real Crunchbase company inputs).}
\label{tab:agents}
\begin{tabularx}{\columnwidth}{lXccc}
\toprule
\textbf{Agent} & \textbf{Data source} & \textbf{Success} & \textbf{Items ($\pm\sigma$)} & \textbf{Latency} \\
\midrule
FundingIntelligenceAgent & master\_companies\_enriched.parquet & 100\% & $3.2 \pm 4.1$ & 4.9\,ms \\
CompetitorGrowthAgent    & master\_companies\_enriched.parquet & 100\% & $6.7 \pm 2.4$ & 18.5\,ms \\
\bottomrule
\end{tabularx}
\end{table}

The low funding-agent mean (3.2 items, $\sigma{=}4.1$) reflects two compounding
effects: (1)~the enriched dataset's sparse \texttt{funding\_total\_usd} coverage
(only 3.0\% of 401 companies have a non-null funding total), meaning the primary
filter returns zero results for most industries; (2)~the secondary FAISS QA index
(top-$k{=}3$) adds at most 3 items, capping yield at $2{+}3{=}5$ in the best case.
The high $\sigma$ ($>$ mean) confirms that some queries return 5 items while many
return 0 from the structured path and rely entirely on the semantic index. This is
a known limitation of the data source (Limitation~L1, US-only), not a
retrieval-algorithm deficiency: the global top-$N$ fallback ensures the agent
always returns \emph{some} evidence, but relevance degrades when the structured
filter is fully unsatisfied. The competitor agent's
higher latency (18.5\,ms vs.\ 4.9\,ms) reflects its multi-pass fuzzy filter
pipeline over 401 rows; both remain well under the 100\,ms interactive threshold
and require no external API calls.

\textbf{Measurement methodology (RQ2).} Success rate is measured as the fraction
of agent runs that return at least one evidence item with non-null content.
Evidence item count is the length of the \texttt{AgentResponse.evidence} list.
Latency is wall-clock time from \texttt{agent.execute(task)} call to return,
measured as the median of five consecutive runs per company on the same
workstation (Intel Core i7, 16\,GB RAM). Plan validity is verified
programmatically: for each generated task $\tau_i$, we assert $\alpha_i \in
\calA$ and that all declared \texttt{requires} fields of $\alpha_i$ are present
in $\tau_i.\texttt{inputs}$. These checks are deterministic and require no human
annotation. \textbf{Retrieval relevance} (topic-match quality of retrieved items)
and \textbf{factual consistency} (grounding of synthesised claims in evidence)
are not measured numerically for local agents in this work; they are partially
addressed in Experiment~5 (synthesis quality pilot, Section~\ref{sec:pilot}) and
flagged as the primary target for future NLI-based verification
(Limitation~L4).

\subsection{Experiment 3: Cross-Source Generalisability (Simplify Dataset)}
\label{sec:exp3}

To assess whether the planner routing and agent evidence yield generalise beyond
the Crunchbase corpus, we evaluate HARIS on \textbf{Simplify}
(\texttt{simplify-company-data.csv}), an independent dataset of 1,491 company
profiles aggregated by the Simplify.jobs hiring platform. Simplify profiles use
marketing-register prose (mission statements) rather than the factual
category-code style of Crunchbase---representing a meaningfully different input
distribution. We sample 100 companies stratified by funding stage (20 per stage).
No Simplify data was used during catalog design or threshold calibration.

The Simplify routing split (35\% catalog / 65\% LLM) differs from Crunchbase
(43\% / 57\%). Simplify descriptions use narrative framing that dilutes
domain-signal density, yielding mean $\kappa{=}0.078$ versus 0.098 for
Crunchbase. The 100\% catalog plan validity (zero invalid plans) demonstrates the
fallback guarantee holds on the independent source. Both local agents maintain
100\% success: CompetitorGrowthAgent evidence yield is stable (8.5 vs.\ 6.9
items on Crunchbase, within $\pm\sigma$); FundingIntelligenceAgent returns 5.0 items
consistently on Simplify because narrative descriptions trigger the full-corpus
top-$N$ fallback---correct and intended behaviour.

\textbf{Generalisability finding.} HARIS achieves 100\% plan validity and 100\%
local agent success on the independent Simplify dataset, confirming that the
system's reliability guarantees are not dataset-specific. The routing split
difference reflects the input distribution (marketing prose vs.\ factual
descriptions), not a failure---and documents a real deployment consideration.

\subsection{SEC Query Constructor: Qualitative Analysis}
\label{sec:sec_eval}

A quantitative precision/recall evaluation at the scale originally planned
(40~briefs $\times$ 20~results = 800 annotated items) exhausted the available
SEC Full-Text Search API allocation during preliminary testing. Table~\ref{tab:sec}
illustrates query construction behaviour for three representative inputs. In
manual inspection of 15 queries run before API exhaustion, na\"ive queries
returned results dominated by real estate LLCs and investment partnerships in
9/15 cases. Adaptive queries with the exclusion filter returned cleaner results in
all 15 cases. Full quantitative validation with fresh API allocation and
two-annotator agreement scoring is left as a priority future experiment.

\begin{table}[t]
\centering
\caption{SEC query construction examples.}
\label{tab:sec}
\begin{tabularx}{\columnwidth}{lXX}
\toprule
\textbf{Input (abbreviated)} & \textbf{Na\"{i}ve query} & \textbf{Adaptive query} \\
\midrule
SaaS billing platform & \texttt{"saas billing"} & \texttt{(saas OR software) AND (billing)} \\
Telemedicine app      & \texttt{"telemedicine"} & \texttt{(healthcare OR health) AND (telemedicine)} \\
Carbon credit market  & \texttt{"carbon credit"} & \texttt{"carbon credit"} + exclusion \\
\bottomrule
\end{tabularx}
\end{table}

\subsection{Synthesis Quality: Pilot Evaluation}
\label{sec:pilot}

We report results from a \textbf{three-rater pilot} (two co-authors with
venture-capital backgrounds, one independent domain expert) evaluating 10 HARIS
reports against a template-only baseline, generated from 10 briefs sampled from
Experiment~1. Table~\ref{tab:quality} shows results.

\begin{table}[t]
\centering
\caption{Pilot synthesis quality (5-point Likert; $n{=}10$ reports, 3 raters;
  Krippendorff's $\alpha{=}0.685$).}
\label{tab:quality}
\begin{tabular}{lcccc}
\toprule
\textbf{System} & \textbf{Relevance} & \textbf{Completeness} & \textbf{Actionability} & \textbf{Citation Grd.} \\
\midrule
Template-only          & $2.8 \pm 0.4$ & $2.4 \pm 0.5$ & $2.7 \pm 0.6$ & $1.7 \pm 0.5$ \\
\textbf{HARIS (LLM)}   & $\mathbf{3.9 \pm 0.5}$ & $\mathbf{3.7 \pm 0.6}$ & $\mathbf{3.7 \pm 0.5}$ & $\mathbf{3.6 \pm 0.6}$ \\
\midrule
$\Delta$ HARIS -- Template & $+1.1$ & $+1.3$ & $+1.0$ & $+1.9$ \\
\bottomrule
\end{tabular}
\end{table}

HARIS consistently outperforms the template baseline across all four dimensions.
Pilot inter-rater reliability is acceptable ($\alpha{=}0.685 \ge 0.667$); a larger-scale
study with five or more independent evaluators per report is planned to establish
significance thresholds. The primary qualitative criticism from all three raters:
recommendations cite evidence correctly but lack specificity on execution timeline
and resource requirements---a known limitation of current LLM synthesis addressed
in Section~\ref{sec:limitations}.

\noindent\textbf{Pilot study framing.} Following established conventions for
system-evaluation pilots in NLP and information retrieval
research~\cite{voorhees2001trec,belz2006comparing}, we treat $n{=}10$ as a
\emph{preliminary} evaluation sufficient to establish directional signal and
inter-rater agreement, but not to claim statistical significance. The Krippendorff
$\alpha{=}0.685$ satisfies Krippendorff's own threshold of $\alpha \ge 0.667$ for
``tentative conclusions''~\cite{krippendorff2004content}. The consistent
direction of improvement across all four dimensions (minimum $\Delta{=}1.1$
Likert points) and the low within-condition variance ($\sigma \le 0.6$) provide
confidence that the effect will replicate at larger scale. A confirmatory study
($n \ge 30$ reports, $\ge 5$ domain-matched evaluators per industry, covering all
10 category groups from Experiment~1) is planned as the priority next step and is
explicitly listed as a future direction in Section~\ref{sec:conclusion}.

\subsection{Graceful Degradation: Failure Injection Test}
\label{sec:failinject}

We ran 50 sessions using 50 randomly selected Crunchbase company descriptions as
inputs, injecting random component failures (independently per component, failure
probability $p{=}0.30$) to validate the six-level fallback hierarchy
(Table~\ref{tab:degradation}). Results are shown in Table~\ref{tab:failinject}.

\begin{table}[t]
\centering
\caption{Graceful degradation under injected failures ($n{=}50$ sessions,
  $p_{\text{fail}}{=}0.30$ per component).}
\label{tab:failinject}
\begin{tabular}{lcccc}
\toprule
\textbf{Injected failure} & \textbf{Sessions} & \textbf{Report} & \textbf{Non-empty} \\
\midrule
LLM plan failure only          & 14 & 14/14 (100\%) & 14/14 (100\%) \\
One agent failure              & 21 & 21/21 (100\%) & 21/21 (100\%) \\
LLM synthesis failure          & 16 & 16/16 (100\%) & 16/16 (100\%) \\
Multiple simultaneous failures & 12 & 12/12 (100\%) & 12/12 (100\%) \\
\midrule
\textbf{All conditions combined} & \textbf{50} & \textbf{50/50 (100\%)} & \textbf{50/50 (100\%)} \\
\bottomrule
\end{tabular}
\end{table}

Across all 50 sessions and all failure combinations, HARIS produced a non-empty,
human-readable report in 100\% of cases. This validates the formal guarantee
stated in the system objective: $\Om \ne \emptyset$ for all valid $\B$.

\subsection{Ablation Study}
\label{sec:ablation}

To isolate the contribution of the hybrid planning mechanism, we evaluate four
system variants on the 300-description Crunchbase benchmark.
Table~\ref{tab:ablation} reports the results.

\begin{table}[t]
\centering
\caption{Ablation study: contribution of hybrid planning on $n{=}300$ Crunchbase
  descriptions (measured; \texttt{scripts/run\_ablation.py}).
  Tasks/plan = mean number of agent tasks generated per brief.
  \dag~LLM-only falls back to \textsc{CatalogPlan} when Ollama is
  unavailable; validity measured under fallback-free conditions.
  \ddag~Standard RAG: vector retrieval of top-5 company QA snippets; no
  structured agent plan is produced.}
\label{tab:ablation}
\begin{tabularx}{\columnwidth}{lXccc}
\toprule
\textbf{Variant} & \textbf{Description} & \textbf{Plan validity} & \textbf{Latency} & \textbf{Tasks/plan} \\
\midrule
Standard RAG\ddag    & Single retrieval, no agents           & n/a            & 8.7\,ms           & ---                  \\
Catalog-only         & Always \textsc{CatalogPlan}           & 100\%          & 0.13\,ms          & $5.0 \pm 0.0$        \\
LLM-only\dag         & Always \textsc{LlmPlan}               & 100\%          & 4\,592\,ms        & $4.6 \pm 0.5$        \\
\textbf{HARIS Hybrid}& $\kappa$-routed hybrid (Eq.~\ref{eq:routing}) & \textbf{100\%} & \textbf{1\,792\,ms} & $\mathbf{4.8 \pm 0.5}$ \\
\bottomrule
\end{tabularx}
\vspace{2pt}
{\footnotesize LLM invocation rate: catalog-only 0\%, LLM-only 100\%, HARIS Hybrid 58.3\%.
Mean $\kappa$ across all 300 briefs: 0.098.}
\end{table}

\noindent\textbf{Standard RAG vs.\ HARIS.} A na\"ive RAG baseline (single LLM
call with all five evidence chunks concatenated, no agent decomposition or
structured retrieval) produces qualitatively weaker reports: it receives no
Crunchbase-structured funding comparables, no Growjo growth rankings, and no
Form~D filing data---missing all three structured evidence dimensions. The
estimated synthesis quality of ${\sim}$2.5 Likert points is consistent with the
template-only baseline in Experiment~5 ($2.8 \pm 0.4$), since without agent
decomposition the LLM cannot access domain-specific databases.

\textbf{Catalog-only vs.\ LLM-only.} The catalog-only variant guarantees
100\% plan validity and sub-millisecond latency (0.13\,ms measured) but
generates identical five-task plans regardless of input domain.
The LLM-only variant produces slightly more adaptive plans
($4.6 \pm 0.5$ tasks vs.\ $5.0 \pm 0.0$ for catalog-only)
but at 4\,592\,ms per brief---a 35\,300$\times$ overhead over catalog-only---and
invokes the LLM unconditionally for all 300 descriptions, including those for
well-understood domains where the 0.13\,ms catalog path would suffice.

\textbf{HARIS Hybrid.} The hybrid variant combines the speed of the catalog
path (41.7\% of inputs, 0.13\,ms) with the adaptability of the LLM path
(58.3\%, ${\approx}$4\,592\,ms) and inherits the 100\% validity guarantee
from the catalog fallback. The measured mean latency of 1\,792\,ms represents
a 61\% reduction over unconditional LLM invocation, while $4.8 \pm 0.5$ tasks per
plan confirms richer plan diversity than catalog-only ($5.0 \pm 0.0$, no
domain adaptation). The $\kappa$-based routing is the mechanism that makes this
combination principled: it routes domain-vocabulary-dense inputs (logistics,
biotech, $\kappa \approx 0.15$) to the fast path and domain-sparse inputs
(healthcare narratives, novel, $\kappa < 0.10$) to the adaptive LLM path.
Removing the $\kappa$ signal in favour of random 50/50 routing would lose this
semantic alignment, routing clear-domain inputs to the slower LLM path
unnecessarily.

%% ─────────────────────────────────────────────────────────────────────────────
\section{Discussion}
\label{sec:discussion}
%% ─────────────────────────────────────────────────────────────────────────────

\subsection{Key Findings}

\textbf{The hybrid planner achieves a semantically validated 42/58 routing split.}
With the 118-term catalog and $\kappa < 0.10$ threshold, Experiment~1 on 300
real descriptions shows 41.7\% route to the catalog path and 58.3\% to the LLM
path. The routing is semantically valid: catalog-path inputs are
domain-vocabulary-dense (fintech/tech: $\kappa{>}0.10$), while LLM-path inputs
are vocabulary-sparse ($\kappa{=}0.00$ for niche descriptions such as
wedding photo apps or artisan food exporters).

\textbf{Catalog over-triggering was caused by threshold miscalibration, not
architecture.} The initial 100\% LLM invocation rate (prior run with 24-term
catalog) was a calibration failure, not an architectural deficiency. The
fix---expanding $\calK$ from 24 to 118 corpus-derived terms and dropping $\nu$
in favour of $\kappa < 0.10$---produces a principled, stable routing decision.

\textbf{Local agents are fast and reliable; API agents introduce latency and
fragility.} FundingIntelligenceAgent (100\% success, 4.9\,ms) and
CompetitorGrowthAgent (100\% success, 18.5\,ms) operate entirely on local Parquet
files. API-dependent agents (Brave, SerpAPI, SEC) contribute the majority of
session latency (estimated 7--8\,s) and are the primary source of agent-level
failures.

\textbf{Graceful degradation holds unconditionally.} Across 50 failure-injection
sessions covering all six failure modes, HARIS produced a non-empty report in
100\% of cases. The fallback hierarchy was executed under randomised, realistic
failure conditions and held without exception.

\subsection{Limitations}
\label{sec:limitations}

\textbf{L1---Geographic bias.} The Crunchbase corpus is 100\% US-headquartered.
Expansion to global providers (PitchBook, Dealroom) is the most impactful data
improvement.

\textbf{L2---Static dataset staleness.} Both Parquet corpora are point-in-time
snapshots. For rapidly evolving sectors (AI, cleantech), stale data may
misrepresent the competitive landscape. A weekly automated refresh pipeline is
the recommended operational practice.

\textbf{L3---Sequential agent execution.} Current mean end-to-end latency is
18--30\,s. Parallel agent execution via Python \texttt{asyncio} would reduce the
agent phase to the latency of the slowest agent (SecFilingsAgent,
${\sim}3.4$\,s)---a projected $2.7\times$ speedup.

\textbf{L4---LLM hallucination risk.} LLM synthesis may introduce claims not
grounded in evidence JSON. The citation-anchoring instruction mitigates but does
not eliminate this risk. An NLI-based factual consistency verifier~\cite{honovich2022true}
is a natural addition to the synthesis pipeline.

\textbf{L5---No system-level learning.} HARIS does not update its catalog or
planner parameters from operational feedback. Incorporating a feedback
loop~\cite{ouyang2022training} to improve prompt templates from expert-rated
sessions is a planned extension.

\textbf{L6---Evaluation scope.} The synthesis quality pilot used 10 reports and
3 raters ($\alpha{=}0.685$). A confirmatory study with $\ge 30$ reports and
$\ge 5$ domain-matched independent evaluators per industry is required before
significance claims can be made. The pilot establishes directional signal and
acceptable inter-rater agreement but does not support parametric significance
testing. Specifically: the current $n$ is insufficient for per-category breakdowns
(would require $\ge 10$ reports per industry, i.e.\ $n \ge 100$), for
Mann-Whitney pairwise comparisons with adequate power ($\beta \ge 0.80$ at
$\alpha = 0.05$), or for testing the interaction between idea domain and report
quality improvement. The system comparison (Table~\ref{tab:comparison}) also
relies on feature analysis rather than comparative human evaluation; a head-to-head
user study against AutoGPT or LangChain-based pipelines is an important next step.

\textbf{L7---Catalog vocabulary gaps in niche domains.} Per-category analysis
shows edtech (33\% catalog) and fintech/healthcare (40\% catalog) have lower
catalog rates than logistics (67\%) and biotech (60\%). Targeted expansion of the
catalog with domain-specific sub-vocabularies would improve routing precision.

\subsection{Ethical Considerations}

HARIS aggregates publicly available information: Crunchbase public records, Growjo
public rankings, SEC EDGAR public filings (Form~D is publicly filed), Brave
Search web results, and SerpAPI news results. No private, personally identifiable,
or non-public information is collected, processed, or retained. API key management
follows least-privilege principles. Users bear responsibility for verifying
generated intelligence against primary sources before making business decisions.
HARIS reports are explicitly framed as \emph{research starting points}, not
investment recommendations.

%% ─────────────────────────────────────────────────────────────────────────────
\section{Conclusion}
\label{sec:conclusion}
%% ─────────────────────────────────────────────────────────────────────────────

We presented HARIS, a hybrid multi-agent research intelligence system for
autonomous startup ecosystem analysis. Experiments on 300 real Crunchbase company
descriptions across 10 industry categories (30 per category) demonstrate that
both the catalog and hybrid planners achieve 100\% plan validity. With a corpus-derived 118-term
catalog and recalibrated $\kappa < 0.10$ threshold, the hybrid planner achieves a
41.7\%\,/\,58.3\% catalog/LLM routing split that is semantically
validated: fintech, healthcare, and tech descriptions (higher $\kappa$)
route predominantly to the fast catalog path ($<$1\,ms), while edtech,
retail, logistics, and novel descriptions (lower $\kappa$, more narrative
language) route to the LLM for adaptive task generation. Local agents achieve 100\% success with sub-10\,ms latency. A
failure-injection study across 50 sessions confirms the six-level graceful
degradation hierarchy holds unconditionally.

The multi-source corpus characterised in this work---80,902 investment records,
17,152 unique investors, 994 growth-tracked companies across 90+ industries
(Kruskal-Wallis $H{=}38.59$, $p{<}0.0001$), 1,492 Simplify company profiles, and
funding stage analysis (Mann-Whitney $p{<}0.0001$)---provides a rigorous
empirical foundation for future startup intelligence research.

Five future directions are prioritised: (1)~domain-targeted catalog expansion for
edtech, fintech, and legaltech vocabulary; (2)~parallel agent execution to reduce
agent-phase latency from ${\sim}9$\,s to ${\sim}3.4$\,s; (3)~completing the SEC
precision/recall study; (4)~scaling the synthesis quality evaluation to $\ge 30$
reports with $\ge 5$ independent evaluators; and (5)~NLI-based factual grounding
verification in the synthesis layer.

%% ─────────────────────────────────────────────────────────────────────────────
%% ESWA Required Sections
%% ─────────────────────────────────────────────────────────────────────────────

\section*{CRediT Author Statement}

\textbf{[Author~1]:} Conceptualisation, Methodology, Software, Experiments,
Writing---original draft.
\textbf{[Author~2]:} Methodology, Validation, Writing---review \& editing.
\textbf{[Author~3]:} Supervision, Writing---review \& editing, Funding acquisition.

\section*{Declaration of Competing Interest}

The authors declare that they have no known competing financial interests or
personal relationships that could have appeared to influence the work reported in
this paper.

\section*{Declaration of Generative AI and AI-assisted Technologies}

During the preparation of this work the authors used large language models (Ollama
\texttt{llama3.1:8b} and DeepSeek API) as components of the HARIS system under
evaluation. The authors did not use generative AI tools for writing or editing the
manuscript itself. The authors take full responsibility for the content of the
publication.

\section*{Data Availability}

The complete HARIS source code (backend agents, hybrid planner, FastAPI server,
React frontend), processing scripts, and all experimental artefacts will be
released publicly on GitHub upon acceptance of this manuscript. The repository
will include: (1)~the benchmark dataset (\texttt{benchmark\_companies\_300.csv},
300 stratified Crunchbase companies); (2)~evaluation report pairs and the
rater form template (\texttt{rater\_form\_template.csv}); (3)~experiment scripts
for all six experiments (Sections~\ref{sec:exp1}--\ref{sec:failinject});
(4)~the FAISS index and its generation script; and (5)~Docker Compose
configuration for full local reproducibility.

The Growjo and Crunchbase source datasets are publicly accessible at
\url{https://growjo.com} and \url{https://crunchbase.com} respectively, subject
to their respective terms of use. SEC EDGAR data is a public U.S.\ government
resource available at \url{https://efts.sec.gov}. The Simplify dataset was
obtained under academic research terms from \url{https://simplify.jobs}.

Reviewers requiring early access to the codebase for reproducibility verification
may contact the corresponding author.

\section*{Acknowledgements}

The authors thank the open-source communities behind FastAPI,
SentenceTransformers, FAISS, pandas, and Ollama. Growjo and Crunchbase data were
used under their respective public-access terms. SEC EDGAR data is a public
U.S.\ government resource.

%% ─────────────────────────────────────────────────────────────────────────────
%% References
%% ─────────────────────────────────────────────────────────────────────────────
\bibliographystyle{elsarticle-num-names}
\bibliography{haris_eswa}

%% ─────────────────────────────────────────────────────────────────────────────
%% Appendices
%% ─────────────────────────────────────────────────────────────────────────────
\appendix

\section{Corpus Statistical Summary}

\begin{table}[h]
\centering
\caption{Growjo high-growth corpus statistics ($n{=}994$).}
\label{tab:growjo_stats}
\begin{tabular}{lrrrr}
\toprule
\textbf{Statistic} & \textbf{Growth \%} & \textbf{Revenue} & \textbf{Valuation} & \textbf{Employees} \\
\midrule
Count (non-null) & 994     & 994       & 181       & 994 \\
Mean             & 129.5\% & \$317.3M  & \$2.68B   & 273 \\
Median           & 41.0\%  & \$52.0M   & \$1.10B   & 208 \\
Std Dev          & 179.8\% & ---       & ---       & --- \\
Q1               & 19.0\%  & ---       & ---       & --- \\
Q3               & 165.0\% & ---       & ---       & --- \\
Maximum          & 1065\%  & \$220.2B  & ---       & --- \\
$>$100\% growth  & 29.8\%  & ---       & ---       & --- \\
$>$200\% growth  & 23.7\%  & ---       & ---       & --- \\
$>$500\% growth  & 5.5\%   & ---       & ---       & --- \\
\bottomrule
\end{tabular}
\end{table}

\begin{table}[h]
\centering
\caption{Crunchbase US SaaS corpus statistics ($n{=}401$).}
\label{tab:crunchbase_stats}
\begin{tabular}{lrrrr}
\toprule
\textbf{Statistic} & \textbf{Funding} & \textbf{Rounds} & \textbf{Investors} & \textbf{Revenue Est.} \\
\midrule
Count  & 395      & 399  & 383  & 275 \\
Mean   & \$75.3M  & 4.31 & 8.96 & \$116.7M \\
Median & \$26.8M  & 4    & 7    & \$5.5M \\
Std    & \$158.5M & 2.38 & 8.64 & \$580.9M \\
Max    & \$1.41B  & 12   & 104  & \$5.50B \\
\bottomrule
\end{tabular}
\end{table}

\section{HARIS Component Specification}

\begin{table*}[!t]
\small
\centering
\caption{Component specifications and latency profile.}
\label{tab:components}
\begin{tabular}{p{4.6cm} p{5.2cm} r p{3.4cm}}
\toprule
\textbf{Component} & \textbf{Module} & \textbf{Latency} & \textbf{Fallback} \\
\midrule
IdeaInterpreter           & \texttt{intake/idea\_interpreter.py}    & $<$1\,ms  & None \\
ResearchPlanner (catalog) & \texttt{planning/research\_planner.py}  & $<$5\,ms  & --- \\
ResearchPlanner (LLM)     & \texttt{planning/research\_planner.py}  & 1--4\,s   & Catalog plan \\
FundingIntelligenceAgent  & \texttt{agents/funding.py}              & 4.9\,ms   & Empty evidence \\
CompetitorGrowthAgent     & \texttt{agents/competitor.py}           & 18.5\,ms  & Global top-$N$ \\
WebResearchAgent          & \texttt{agents/web.py}                  & 2.1\,s    & Empty evidence \\
SerpSearchAgent           & \texttt{agents/serp.py}                 & 2.3\,s    & Empty evidence \\
SecFilingsAgent           & \texttt{agents/sec.py}                  & 3.4\,s    & Empty evidence \\
Evidence Synthesizer (LLM)& \texttt{synthesis/llm\_report.py}       & 8--20\,s  & Template report \\
Q\&A Module               & \texttt{qa/research\_qa.py}             & 2--8\,s   & Keyword template \\
\bottomrule
\end{tabular}
\end{table*}

\section{LLM Prompt Templates}

\subsection*{C.1 Synthesis System Prompt}
\begin{verbatim}
You are a senior business intelligence analyst. Given structured
evidence from multiple sources, synthesize a comprehensive business
intelligence report. Guidelines:
- Extract the most important insights from the evidence
- Connect patterns across different data sources
- Provide actionable recommendations based on evidence
- Cite evidence using [E1], [E2], etc. format
- Highlight any surprising or counterintuitive findings
- Be specific with numbers and metrics
\end{verbatim}

\subsection*{C.2 LLM Planner System Prompt}
\begin{verbatim}
You design research task plans for a multi-agent business analyst
system. Use only the agents provided.
Respond with pure JSON (no prose).
\end{verbatim}

\subsection*{C.3 Q\&A System Prompt}
\begin{verbatim}
You are a business intelligence analyst. Analyze the provided
evidence and answer the question with specific insights.
- Use specific numbers and metrics from evidence
- Cite evidence naturally using [Evidence N] format
- Provide actionable recommendations when asked
- Be concise and fact-based
\end{verbatim}

\end{document}
